%% --------------------------------------------------------------------------
\section{Property Graph Queries}
\label{sec:eval:property}

We integrate \project{} into the AsterDB evaluation artifact~\cite{moAsterEnhancingLSMstructures2025}, which provides a reproducible Docker-based benchmark suite that compares seven graph storage systems: Neo4j, ArangoDB, JanusGraph, OrientDB, PostgreSQL (via SQLG), NebulaGraph, and AsterDB on structural and property graph workloads.
The AsterDB artifact uses Apache TinkerPop's Gremlin traversal language~\cite{apachetinkerpop} as a uniform query interface: each graph database implements a TinkerPop \emph{provider} that translates Gremlin traversals into its native storage operations, enabling an apples-to-apples comparison across backends.
\project{} bypasses this layer entirely, running benchmarks through its native C++ API; we quantify the overhead this avoids in~\autoref{sec:eval:property-ldbc}.
We evaluate \project{}'s property graph capabilities in two complementary experiments.
First, we run graph analytics benchmarks on a rather small structural graph (cit-Patents).
Second, we design a query benchmark using 14 queries derived from the LDBC Social Network Benchmark (SNB)~\cite{erling2015ldbc} at scale factor~3 (${\sim}800$\,K vertices, ${\sim}3.5$\,M edges across 4~entity types and 7~relation types).
We compare four \project{} configurations (Adjacency List and EdgeKey, each with embedded and columnar property storage; \autoref{arch:prop-embedded}, \autoref{arch:prop-columnar}) to Neo4j (Cypher), ArangoDB (AQL), and NeuG~\cite{NeuG2025Alibaba} (Cypher, embedded mmap-based CSR).

We exclude AsterDB from the LDBC comparison: although its RocksDB backend persists properties, the Gremlin bridge reads vertex properties from an in-memory map that is never populated from storage, so property lookups return null; additionally, all edges are stored under a single default label, preventing label-filtered traversals.
% AsterDB's property write path and indexed \code{has()} lookups are fully wired, so it participates in the CRUD microbenchmarks.

\section{Graph algorithms.}
\label{sec:eval:property-graphalytics}
\autoref{tab:graph-algorithms} reports algorithm runtimes on cit-Patents (3.77\,M vertices, 16.5\,M edges).
PageRank (PR), community detection (CDLP), and weakly-connected components (WCC) are full-graph iterative algorithms; BFS and shortest path (SP) are single-pair traversals matching the Gremlin baselines.
\project{} Adjacency~List completes PageRank in 0.3\,s, more than a $1000\times$ faster than Neo4j (429\,s) and nearly $8000\times$ faster than OrientDB (2{,}560\,s).
WCC finishes in 0.3\,s vs.\ Neo4j's 998\,s ($\approx 3000\times$).
These systems are designed for OLTP queries and assume that when analytics are necessary, they are part of an ETL workflow, exporting the structural graph for external processing using a specialized graph processing system.
\project{} lets us bypass this costly ETL process without jeopardizing query performance.
ArangoDB and JanusGraph exceed the 2-hour timeout on all iterative algorithms but complete the single-pair traversals quickly (BFS in 28 to 97\,ms, SP in 35 to 110\,ms) where their native shortest-path implementations apply.
PostgreSQL also times out on all iterative algorithms and takes $\approx$11\,s for the single-pair traversals.
\project{} runs parallel C++ algorithm implementations directly on the \wt{} B-tree; graph databases execute these benchmarks using the TinkerPop abstraction layer. %graph databases execute interpreted procedures through their query engines, adding per-iteration overhead that compounds across the thousands of supersteps these algorithms require.
To quantify the per-call overhead of the Gremlin/TinkerPop query layer, we decompose Neo4j's PageRank execution into three variants on cit-patents:
(1)~\emph{no-op}, which measures the round-trip latency between the TinkerPop client and the database server with no graph access;
(2)~\emph{vertex-lookup}, which additionally resolves a vertex by its index; and
(3)~the full PageRank step, which traverses the vertex's neighbors.
The per-call costs are 4.34 µs, 6.41 µs, and 12.75 µs respectively, revealing that 50\% of Neo4j's PageRank time is spent on Gremlin framework overhead (34\%) and vertex index resolution (16\%) rather than actual edge traversal.
Even if this overhead were entirely eliminated, \project{} remains $\approx700\times$ faster.

%%%%%%%%%%%%%%%%%%%%%%%%%%%%%%%%%%%%%%%%%%%%%%%%%%%%%%%%%%%%%%%%%%%%%%%%%%%%%%%%
\begin{table}[h!]
  \centering
  \footnotesize
  \setlength{\tabcolsep}{2pt}
  \begin{tabular*}{\columnwidth}{@{\extracolsep{\fill}}lrrrrr@{}}
    \toprule
    System & PR & CDLP & WCC & SP & BFS \\
    \midrule
    FG:AdjList & \textbf{0.32} & \textbf{1.67} & \textbf{0.32} & \textbf{0.02} & 0.23 \\
    FG:EdgeKey & 9.60 & 3.24 & 1.03 & 0.02 & 0.29 \\
    AsterDB & 908.53 & 626.14 & 2{,}219.91 & 0.34 & 0.06 \\
    Neo4j & 428.73 & 652.04 & 998.41 & 0.41 & 0.13 \\
    ArangoDB & t/o & t/o & t/o & 0.03 & \textbf{0.03} \\
    OrientDB & 2{,}559.81 & 4{,}750.27 & t/o & 0.22 & 0.06 \\
    PostgreSQL & t/o & t/o & t/o & 11.26 & 11.45 \\
    JanusGraph & t/o & t/o & t/o & 0.11 & 0.10 \\
    \bottomrule
  \end{tabular*}
  \caption{Graph algorithm runtimes (s) on cit-Patents. PR, CDLP, WCC are full-graph iterative; BFS, SP are single-pair. Timout ~=~2\,h.\PM{check the numbers. Also match the prose.}}
  \label{tab:graph-algorithms}
\end{table}
%%%%%%%%%%%%%%%%%%%%%%%%%%%%%%%%%%%%%%%%%%%%%%%%%%%%%%%%%%%%%%%%%%%%%%%%%%%%%%%%



%%%%%%%%%%%%%%%%%%%%%%%%%%%%%%%%%%%%%%%%%%%%%%%%%%%%%%%%%%%%%%%%%%%%
\begin{table}[t]
  \centering
  \footnotesize
  \setlength{\tabcolsep}{3pt}
  \begin{tabular*}{\columnwidth}{@{\extracolsep{\fill}}ll rrrrr@{}}
    \toprule
    & & \multicolumn{2}{c}{\project{} AdjList} & & & \\
    \cmidrule(lr){3-4}
    Cat. & Query & Emb & Col & Neo4j & Arango & NeuG \\
    \midrule
    \multirow{2}{*}{Point (\textmu{}s)}
    & R1   & \textbf{7.28}  & 14.2  & 28.4  & 19.8  & 802 \\
    & X1   & \textbf{6.05}  & 7.63  & 22.1  & 15.2  & 544 \\
    \midrule
    \multirow{3}{*}{Aggr (\textmu{}s)}
    & A2   & \textbf{12.8}  & 16.9  & 66.2  & 38.4  & 805 \\
    & A3   & \textbf{13.8}  & 22.2  & 82.3  & 64.2  & 850 \\
    & X5   & 10.2  & \textbf{7.68}  & 55.8  & 29.1  & 866 \\
    \midrule
    \multirow{3}{*}{Multi (ms)}
    & X3   & \textbf{0.24}  & 1.9  & 13.1  & 13.3  & 1.1 \\
    & IC-3   & 29.0  & 29.0  & 7.6  & 46.0  & \textbf{3.8} \\
    & IC-9   & \textbf{1.5}  & 17.7  & 65.4  & 188.8  & 29.8 \\
    \midrule
    \multirow{2}{*}{Scan (s)}
    & BI-1   & ---  & \textbf{0.40}  & 1.78  & 6.67  & 0.89 \\
    & BI-12   & ---  & 1.24  & 2.66  & 1.51  & \textbf{0.09} \\
    \midrule
    \multirow{4}{*}{Write (\textmu{}s)}
    & W1   & 110  & 413  & 26.2  & \textbf{21.3}  & 190{,}420 \\
    & W2   & 165  & 454  & 30.9  & \textbf{26.9}  & 192{,}304 \\
    & W3   & 170  & 436  & 68.0  & \textbf{60.7}  & --- \\
    & X4   & \textbf{11.3}  & 117  & 37.4  & 20.4  & --- \\
    \bottomrule
  \end{tabular*}
  \caption{LDBC SNB query p50 latency, scale factor~3.
    Units: \textmu{}s for point, aggregate, and write queries;
    ms for multi-hop queries; seconds for scan queries.
    Neo4j and ArangoDB point/aggregate queries measured via batched execution
    (1000 queries per request) to amortise per-request overhead.
    Emb~=~embedded properties, Col~=~columnar.
    Best per row in \textbf{bold}; ---~=~not applicable.}
  \label{tab:ldbc-queries}
\end{table}

%%%%%%%%%%%%%%%%%%%%%%%%%%%%%%%%%%%%%%%%%%%%%%%%%%%%%%%%%%%%%%%%%%%%
\section{LDBC SNB Query Performance}
\label{sec:eval:property-ldbc}
% We run 10~queries derived from the LDBC SNB specification, organized into four categories (\autoref{tab:ldbc-queries}).
% \emph{Point queries} (R1, X1) fetch all properties of a single vertex.
% \emph{Aggregates} (A2, A3, X5) count edges matching a date-range predicate.
% \emph{Multi-hop traversals} (X3, IC-3, IC-9) combine 2-hop structural traversal with property filters.
% \emph{BI scans} (BI-1, BI-12) aggregate over all vertices of a type.
% We also measure write operations (W1, W2).

We benchmark 14 queries organized into five categories
(\autoref{tab:ldbc-queries}): point reads, aggregates, multi-hop traversals,
BI scans, and writes.
Five queries (IC-3, IC-9, X3, BI-1, BI-12) are derived from the LDBC SNB
Interactive and Business Intelligence workloads; we note simplifications below.
The remaining read queries (R1, X1, A2, A3, X5) are targeted microbenchmarks
that isolate individual access patterns (single-vertex property fetches and
edge-count aggregates with temporal predicates) that the LDBC complex queries
bundle with multi-hop traversal, making it difficult to attribute latency to
the property layer alone.
% \autoref{tab:query-desc} summarizes each query.

Following the LDBC SNB parameter curation methodology, we select query
parameters that produce varying result-set sizes to exercise both best- and
worst-case fan-out.
We sample 90~Person IDs via degree-stratified tercile sampling: all Persons are
sorted by \textsc{knows} out-degree, split into three equal-sized terciles
(low, medium, high degree), and 30~IDs are drawn uniformly at random from each
tercile (deterministic seed for reproducibility).
Post parameters are generated similarly, stratified by incoming
\textsc{likes} count.

Neo4j and ArangoDB operate as client-server systems, so every individual query pays for network round-trip, query parsing, and result serialisation.
Single-shot read query measurements reported Neo4j R1 at $1.95$\,ms, but $\approx$98\% of that was per-request network and query-parse overhead: Neo4j's no-op round-trip alone costs $\approx$1.5\,ms.
NeuG and \project{}, by contrast, are embedded stores with no network layer.
To avoid penalising Neo4j and ArangoDB for an architectural choice orthogonal to query-engine performance, we batch 1,000 parameterised queries in a single \texttt{UNWIND} request (Cypher) or \texttt{FOR-IN} loop (AQL) and divide by the batch size, amortising per-request overhead to $<$3\,$\mu$s.
NeuG requires no such batching: it is embedded and its query plan cache (keyed by query string) eliminates repeated planning.
Each query is run 100~times after warmup, cycling through the 90~parameter
values; we report the median (p50) latency.
\autoref{tab:ldbc-queries} reports the results.

\textbf{Point queries and aggregates.}
Point queries (R1, X1) fetch all properties of a single vertex.
R1 returns a Person's 7~scalar properties plus multi-valued email and language
lists; X1 returns a Post's creation date, length, and content.
Aggregates (A2, A3, X5) scan a vertex's edge list and count edges
whose \texttt{creationDate} falls within a supplied date range.
A2 counts outgoing \textsc{knows} edges, A3 counts outgoing \textsc{likes}
edges (to Posts), and X5 counts incoming \textsc{likes} edges on a given Post.
\project{} Adjacency~List with embedded storage dominates all point and aggregate queries.
R1 completes in $7$\,$\mu$s on \project{}, $3.7\times$ faster than Neo4j ($27$\,$\mu$s), $2.7\times$ faster than ArangoDB ($20$\,$\mu$s), and $110\times$ faster than NeuG ($802$\,$\mu$s).

A potential concern with a direct C++ API is developer effort compared to declarative query languages such as Cypher and Gremlin.
In practice, graph queries are rarely written ad hoc; they are embedded in application logic and invoked programmatically~\cite{10.1145/3357766.3359541,cheung2013optimizing}.
Modern code-generation tools (e.g., LLM-based assistants) further lower this barrier, making it straightforward to generate correct API calls to \project{}'s interface from a high-level description \cite{lyu2026text2gqlbenchtextgraphquery}.

\textbf{Multi-hop traversals.}
Multi-hop queries (X3, IC-3, IC-9) combine structural traversal
with property filters.
X3 implements a simplified variant of LDBC IC-2: given a Person, it traverses
outgoing \textsc{knows} edges to find friends, hops to their Posts via reverse
\textsc{hasCreator} edges, filters by \texttt{creationDate~$<$~cutoff}, and
returns the 10~most recent posts.
Unlike IC-2, X3 retrieves only Posts (not Comments).
IC-3 follows LDBC IC-3: a 2-hop \textsc{knows} expansion to find
friends-of-friends located in a given country, ranked by common-friend count
(top~10).
IC-9 follows LDBC IC-9: friends' messages (Posts and Comments) created before a
cutoff date, ranked by recency (top~20).
Our IC-9 traverses only direct friends (1-hop), whereas the LDBC specification
includes friends-of-friends (2-hop).
%
X3 favors \project{} at $0.24$\,ms, $4.7\times$ faster than NeuG ($1.14$\,ms) and $56\times$ faster than Neo4j ($13.5$\,ms).
IC-9 also favors \project{} ($1.47$\,ms) over Neo4j ($64.5$\,ms, $44\times$) and ArangoDB ($188.8$\,ms, $129\times$).
However, IC-3 reverses the ranking: NeuG completes in $3.84$\,ms, Neo4j in $8.3$\,ms, and \project{} in $29.0$\,ms.
IC-3 fully expands the 2-hop \code{knows} neighborhood, which is large for well-connected persons.
NeuG's in-memory CSR iterates each vertex's neighbor list in a single sequential memory scan per hop, while \project{} must issue a \wt{} B-tree cursor seek and page-level scan for each neighbor list, paying the cost of B-tree traversal at every hop.

\textbf{BI scans.}
BI-1 and BI-12 are analytical queries over all vertices of a
type.
BI-1 follows LDBC BI-1: scan all Post vertices, bucket by creation year and
message-length category (short, medium, long), and return per-bucket counts and
length sums.
BI-12 follows LDBC BI-12: scan all Posts, traverse \textsc{hasCreator} to
identify the author, and return the top-10 creators by post count.
%
BI queries require columnar storage: scanning all Posts in embedded mode would deserialize the full property blob per vertex, negating the scan advantage.
In columnar mode, BI-1 takes $404$\,ms, $4.03\times$ faster than Neo4j ($1{,}629$\,ms) and $16.5\times$ faster than ArangoDB ($6{,}666$\,ms); NeuG takes $892$\,ms.
The columnar layout reads only the temporal and length column groups, skipping content columns entirely.
For BI-12, NeuG excels at $93$\,ms, $13.4\times$ faster than \project{} ($1{,}245$\,ms).
BI-12 adds an edge hop per Post; NeuG's CSR iterates these hops sequentially in memory, while \project{} pays a B-tree cursor seek per hop.

\textbf{Write operations.}
W1--W3 and X4 insert vertices and edges with properties.
W1 inserts a Person vertex with all scalar properties and an
\textsc{isLocatedIn} edge (LDBC Insert~1).
W2 inserts a bidirectional \textsc{knows} edge between two Persons (LDBC
Insert~4).
W3 inserts a Post vertex with a \textsc{hasCreator} edge to its author (LDBC
Insert~6).
X4 inserts a \textsc{likes} edge from a Person to a Post (LDBC Insert~5).
%
ArangoDB and Neo4j dominate W1--W3 ($21$--$68$\,$\mu$s) because their write paths are optimized for single-document upserts through mature, heavily tuned storage engines.
\project{} takes $110$--$170$\,$\mu$s in embedded mode; each write opens a WiredTiger transaction, performs a B-tree cursor insert, and commits, overhead that is amortized in bulk ingest but visible at single-operation granularity.
X4 (single edge insert with no vertex creation) favors \project{} at $11$\,$\mu$s, where the cost reduces to a single cursor insert.
NeuG is $1{,}000\times$ slower on all writes ($190$\,ms) because each mutation parses and plans a full Cypher \texttt{CREATE} statement.

\section{Embedded vs.\ Columnar Trade-off}

The two property storage modes offer a clear, query-dependent trade-off (\autoref{tab:ldbc-queries}, Emb vs.\ Col columns).
For point queries and single-hop reads, both modes perform similarly (R1: $7$\,$\mu$s embedded, $14$\,$\mu$s columnar), because the dominant cost is the structural cursor seek, not the property read.
For multi-hop traversals that access properties at each hop, embedded mode wins decisively: IC-9 takes $1.47$\,ms embedded vs.\ $17.7$\,ms columnar ($12\times$); IC-5 (not shown) takes $2.65$\,ms vs.\ $15.2$\,ms ($5.7\times$).
The overhead stems from columnar mode's per-hop cost: at each traversal step, the query opens a cursor on the target column group's B-tree, while embedded mode reads properties inline from the already-fetched vertex or edge value.

Conversely, BI queries that scan a single column benefit from columnar storage.
The column-group B-tree contains only the accessed columns (e.g., \code{creationDate} and \code{length} for BI-1), achieving high row density per page and efficient sequential scans.
Embedded mode is impractical for these queries, because every row's full property blob must be deserialized to extract one field.

This trade-off suggests a deployment strategy: use embedded mode for OLTP-style queries (point lookups, narrow multi-hop traversals) and columnar mode for analytical scans over single property columns.
In practice, the choice can be made per edge or vertex label, since each label's property table is independent.
